\section{Formal Definitions}\label{sec:defs}
\heading{Single-server pre-processing PIR.}
We first define a single-server pre-processing PIR scheme. 
A single-server pre-processing PIR scheme consists of two stateful algorithms: the client and the server. The scheme consists of the two following phases.
\begin{enumerate}
    \item {\bf Pre-processing :} The pre-processing is run only once at the beginning. The client receives no input, while the server receives a database $\DB\in \bits^n$ as input. The client and server interact and the client may store some information in its local storage. We refer to this information as ``hints''.
    \item {\bf Queries:} This phase is repeated for every index $x$ of the $\DB$ that the client wants to read. For every query, the client sends a single message to the server, and the server responds with a single message. The client then performs some computation and outputs an answer $\beta$.
\end{enumerate}

\paragraph{Correctness.} Given a database $\DB$ with entries indexed by $0,1,\ldots,n-1$, correctness entails that the query for an index $x\in \{0,1.\ldots,n-1\}$ by the client in the query phase, results in an answer $\DB[x]$ (the $x$-th bit of $\DB$) output by the client. 
%We actually consider a slightly relaxed notion of correctness where the scheme is allowed to be incorrect with a small probability.
Formally, correctness requires that for any 
%statistical security parameter $\kappa \in \mathbb{N}$, 
security parameter $\lambda \in \mathbb{N}$, 
for any $n,q$ which are polynomially bounded in $\lambda$, 
there exists a negligible function $\negl$ such that for any database 
$\DB\in \bits^n$, for any sequence of queries $x_1,x_2,\ldots,x_q\in \{0,1,\ldots,n-1\}$, an honest execution of the PIR scheme with $\DB$ and queries $x_1,x_2,\ldots,x_q$ returns all correct answers with probability 
at least $1-\negl(\lambda)$.

\paragraph{Privacy.} Privacy of a PIR scheme entails that for any index $x$ queried by the client to the server, 
the view of the server must not leak information about the query $x$. 
Formally, we define the privacy of a single-server PIR scheme as follows.

A single-server PIR scheme satisfies privacy if and only if there exists a probabilistic polynomial-time simulator $\Sim(1^\lambda,n)$ such that for any probabilistic polynomial-time adversary $\advA$ acting as the server, any polynomially bounded $n$ and $q$, any database $\DB\in \bits^n$, $\advA$'s
views in the following two experiments are computationally indistinguishable:
\begin{itemize}
    \item $\mathsf{Real}$:  an honest client interacts with $\advA(1^\lambda,n,\DB)$ who acts as the server and may arbitrarily deviate from the prescribed protocol. In every query step $t\in [q]$, $\advA$ may adaptively choose the next query $x_t\in \{0,1,\ldots,n-1\}$ for the client, and the client is invoked with $x_t$ as input.
    \item $\mathsf{Ideal}$: the simulated client $\Sim(1^\lambda,n)$ interacts with $\advA(1^\lambda,n,\DB)$ who acts as the server and may arbitrarily deviate from the prescribed protocol. In every  query step $t\in [q]$, $\advA$ may adaptively choose the next query $x_t\in \{0,1,\ldots,n-1\}$ for the client, and the client is invoked without $x_t$ as input.
\end{itemize}

\paragraph{Two-server pre-processing PIR.}
In the two-server setting, there are two non-colluding servers, and the client
may interact with both servers in both the preprocessing and query phases. 
The two servers do not interact with each other.

Correctness is defined in the same way as the single-server setting. 
For privacy, we want the definition of the single-server setting
to hold for each individual server. 
%Privacy for a two-server pre-processing PIR scheme are analogously defined, with the change that there are two servers now instead of one. For the privacy definition, the adversary controls one of the two servers.

\paragraph{Additional notation.}
In the formal sections later, for clarity we  
distinguish
between a statistical security parameter denoted $\kappa$
and a computational security parameter denoted $\lambda$.
